\section{Supplementary Results and Research Leads}
\label{app:supplementary-register-05}

This appendix keeps potentially relevant results from the research
archive attached to a mathematical home without enlarging the main
theorem spine.  This register is published separately from the six reader manuscripts.  Inclusion here is deliberately permissive.  Each entry
states its evidence boundary; entries labelled as open, contextual, or
evidence-only are not used as theorems elsewhere in the paper.  Earlier
formulations known to be false are replaced by their current correction.
The computational supplement and its \texttt{COMPUTATION.md} index give
the available program-level artifact map; an entry here does not turn a
shared-code check into an independent verification.
\begingroup\raggedright\emergencystretch=3em

\subsection{Additional results and constructions}

\phantomsection
\label{supp-note-05-001}
\paragraph{An explicit degree-three counterexample in 11 variables.}
\emph{Status.} Proof offered; independent review is not recorded.

The cited map from complex affine 11-space to itself has total degree 3, 52 nonzero monomial terms, constant Jacobian determinant -2, and three displayed distinct rational inputs with one common image.

\emph{Credit.} Credit recorded to ChatGPT (research assistance); Spacerat (computation).

\phantomsection
\label{supp-note-05-002}
\paragraph{An explicit cubic-homogeneous counterexample in 24 variables.}
\emph{Status.} Proof offered; independent review is not recorded.

The cited explicit map has the form G(U)=U+H(U) over Q\textasciicircum{}24, with every nonzero component of H homogeneous cubic, determinant 1, 54 nonzero cubic monomials, and a displayed collision of two distinct rational points.

\emph{Credit.} Credit recorded to William Thompson (computation).

\phantomsection
\label{supp-note-05-003}
\paragraph{Starting from the public 11D counterexample, rank-sensitive cubic...}
\emph{Status.} Recorded result or structural lead; it is not promoted beyond the evidence described here and in the project archive.

Starting from the public 11D counterexample, rank-sensitive cubic suspension yields an explicit 19D cubic-homogeneous Keller counterexample with collision.

\phantomsection
\label{supp-note-05-004}
\paragraph{An explicit square-zero Druzkowski pairing of the 19D cubic gives...}
\emph{Status.} Recorded result or structural lead; it is not promoted beyond the evidence described here and in the project archive.

An explicit square-zero Druzkowski pairing of the 19D cubic gives a 135D Druzkowski counterexample with collision.

\phantomsection
\label{supp-note-05-005}
\paragraph{The earlier 26D cubic-homogeneous and 158D Druzkowski constructions...}
\emph{Status.} Recorded result or structural lead; it is not promoted beyond the evidence described here and in the project archive.

The earlier 26D cubic-homogeneous and 158D Druzkowski constructions remain valid explicit examples but are superseded as upper bounds by the later 19D and 135D examples.

\phantomsection
\label{supp-note-05-006}
\paragraph{For an n-dimensional degree-at-most-three map whose quadratic span...}
\emph{Status.} Proof offered; independent review is not recorded.

For an n-dimensional degree-at-most-three map whose quadratic span has rank r, a block-determinant suspension produces a cubic-homogeneous map in dimension n+r+1.

\phantomsection
\label{supp-note-05-007}
\paragraph{For the 19D example, the generic Jacobian perturbation JH has...}
\emph{Status.} Recorded result or structural lead; it is not promoted beyond the evidence described here and in the project archive.

For the 19D example, the generic Jacobian perturbation JH has Jordan type (18, 1); nonzero scalar fibers are conjugate and the zero fiber is an automorphism degeneration.

\phantomsection
\label{supp-note-05-008}
\paragraph{The explicit constructions give the upper bound N\_min <= 19...}
\emph{Status.} Proof offered; independent review is not recorded.

The explicit constructions give the upper bound N\_min <= 19 for cubic-homogeneous counterexamples; the conversation uses the elementary no-collinear-collision argument for the lower bound N\_min >= 5.

\phantomsection
\label{supp-note-05-009}
\paragraph{The five-attacks package constructs a 38D Hessian quartic with...}
\emph{Status.} Recorded result or structural lead; it is not promoted beyond the evidence described here and in the project archive.

The five-attacks package constructs a 38D Hessian quartic with Hessian Jordan type (35, 2, 1) and derives exact central-binomial inverse-ray coefficients with infinitely many nonzero Laplacian terms in the stated parity class.

\emph{Credit.} Credit recorded to ChatGPT (AI-assisted derivation and drafting).

\phantomsection
\label{supp-note-05-010}
\paragraph{The explicit pairing dictionary gives a 110D Druzkowski...}
\emph{Status.} Recorded result or structural lead; it is not promoted beyond the evidence described here and in the project archive.

The explicit pairing dictionary gives a 110D Druzkowski construction, while the derivative-space zero-block and equality-case argument give the model-relative bound 52 <= pairing rank <= 110 for the fixed 19D tensor.

\phantomsection
\label{supp-note-05-011}
\paragraph{The unique scalar-monomial quadratic shift lowering the cubic jet...}
\emph{Status.} Recorded result or structural lead; it is not promoted beyond the evidence described here and in the project archive.

The unique scalar-monomial quadratic shift lowering the cubic jet from rank seven to six is P2=-d\textasciicircum{}2 e\_a, and its quartic error is detected by an exact 13-term functional outside the cubic homological image; the natural quadratic target cleanup also fails.

\phantomsection
\label{supp-note-05-012}
\paragraph{An exact row-defect theorem holds in the two-matrix square-zero...}
\emph{Status.} Recorded result or structural lead; it is not promoted beyond the evidence described here and in the project archive.

An exact row-defect theorem holds in the two-matrix square-zero family; the accompanying finite-field search found no better example in the tested range.

\phantomsection
\label{supp-note-05-013}
\paragraph{A second counterexample \(G:\mathbb C^4\to\mathbb C^4\) has...}
\emph{Status.} Recorded result or structural lead; it is not promoted beyond the evidence described here and in the project archive.

A second counterexample \(G:\mathbb C^4\to\mathbb C^4\) has determinant \(4\) and generic degree \(8\).

\phantomsection
\label{supp-note-05-014}
\paragraph{A straightforward explicit reduction produced a cubic-homogeneous...}
\emph{Status.} Recorded result or structural lead; it is not promoted beyond the evidence described here and in the project archive.

A straightforward explicit reduction produced a cubic-homogeneous map in \(79\) variables and a Dru\.{z}kowski cubic-linear map in \(426\) variables.

\phantomsection
\label{supp-note-05-015}
\paragraph{\(\mathrm{MC}(SU(79))\) is false.}
\emph{Status.} Exact or computer-assisted certificate offered; see the computation index for its scope and independence boundary.

\(\mathrm{MC}(SU(79))\) is false.

\phantomsection
\label{supp-note-05-016}
\paragraph{The new 24-variable map improves this to \(\mathrm{MC}(SU(N))\)...}
\emph{Status.} Recorded result or structural lead; it is not promoted beyond the evidence described here and in the project archive.

The new 24-variable map improves this to \(\mathrm{MC}(SU(N))\) false for all \(N\ge24\).

\phantomsection
\label{supp-note-05-017}
\paragraph{The rank-sensitive cubic suspension is triangularly left-right...}
\emph{Status.} Proof offered; independent review is not recorded.

The rank-sensitive cubic suspension is triangularly left-right equivalent to (X, v, t) mapped by (t\textasciicircum{}\{-1\}K(tX), v, t); it preserves generic degree and the geometric Galois closure after adjoining transcendental variables.

\emph{Credit.} Credit recorded to ChatGPT (AI-assisted derivation and drafting).

\phantomsection
\label{supp-note-05-018}
\paragraph{For any Keller map F, the symmetric double (F(x), JF(x)\textasciicircum{}T y)...}
\emph{Status.} Proof offered; independent review is not recorded.

For any Keller map F, the symmetric double (F(x), JF(x)\textasciicircum{}T y) is polynomially right-equivalent to F times the identity, so the 38-variable gradient descendant has the same generic degree and monodromy as the 19-variable cubic map.

\emph{Credit.} Credit recorded to ChatGPT (AI-assisted derivation and drafting).

\phantomsection
\label{supp-note-05-019}
\paragraph{A full-rank square-zero Gorni-Zampieri pairing H=B(Dz)\textasciicircum{}\{*3\}, BD=0...}
\emph{Status.} Proof offered; independent review is not recorded.

A full-rank square-zero Gorni-Zampieri pairing H=B(Dz)\textasciicircum{}\{*3\}, BD=0, is triangularly right-equivalent to the paired cubic map times an identity factor; the 110-variable Druzkowski map is therefore a stable presentation of the 19-variable cubic map.

\emph{Credit.} Credit recorded to ChatGPT (AI-assisted derivation and drafting).

\phantomsection
\label{supp-note-05-020}
\paragraph{The original three-variable map has the displayed invariant quotient...}
\emph{Status.} Recorded result or structural lead; it is not promoted beyond the evidence described here and in the project archive.

The original three-variable map has the displayed invariant quotient cubic theta\textasciicircum{}3-2theta\textasciicircum{}2+Btheta-2A, whose irreducibility and nonsquare discriminant give generic Galois group S3; the descendant ladder presents this same three-sheeted cover.

\emph{Credit.} Credit recorded to ChatGPT (AI-assisted derivation and drafting).

\phantomsection
\label{supp-note-05-021}
\paragraph{For the fixed 38-variable Hessian quartic, Delta\textasciicircum{}m(Q\textasciicircum{}\{m+1\}) is...}
\emph{Status.} Exact or computer-assisted certificate offered; see the computation index for its scope and independence boundary.

For the fixed 38-variable Hessian quartic, Delta\textasciicircum{}m(Q\textasciicircum{}\{m+1\}) is nonzero for every m>=1; an explicit two-parameter inverse ray removes the earlier parity gap and ties every nonzero coefficient to a sheet escaping at the discriminant.

\emph{Credit.} Credit recorded to ChatGPT (AI-assisted derivation and drafting).

\phantomsection
\label{supp-note-05-022}
\paragraph{Although every JH(z) is nilpotent, the product JH(e\_d)JH(e\_t) has...}
\emph{Status.} Exact or computer-assisted certificate offered; see the computation index for its scope and independence boundary.

Although every JH(z) is nilpotent, the product JH(e\_d)JH(e\_t) has characteristic polynomial lambda\textasciicircum{}18(lambda+3); hence the polynomial-matrix family is not strongly nilpotent and has no common constant triangularizing basis.

\emph{Credit.} Credit recorded to ChatGPT (AI-assisted derivation and drafting).

\phantomsection
\label{supp-note-05-023}
\paragraph{The equivariant cubic rank-six compression condition has a...}
\emph{Status.} Recorded result or structural lead; it is not promoted beyond the evidence described here and in the project archive.

The equivariant cubic rank-six compression condition has a 20-dimensional affine slice, and the quartic obstruction functional is identically one on that slice; modulo two target-row directions this is the natural transverse compression space.

\emph{Credit.} Credit recorded to ChatGPT (AI-assisted derivation and drafting).

\phantomsection
\label{supp-note-05-024}
\paragraph{The n+r+1 cubic suspension is coordinate-free on V plus the...}
\emph{Status.} Proof offered; independent review is not recorded.

The n+r+1 cubic suspension is coordinate-free on V plus the dual of the output-mode span plus one scalar, and r=rank(C\textasciicircum{}flat) is the minimum possible auxiliary dimension for any factorization of the cubic tensor through an intermediate vector space.

\emph{Credit.} Credit recorded to ChatGPT (AI-assisted derivation and drafting).

\phantomsection
\label{supp-note-05-025}
\paragraph{Within the fixed 38-dimensional symmetric double, there is no proper...}
\emph{Status.} Recorded result or structural lead; it is not promoted beyond the evidence described here and in the project archive.

Within the fixed 38-dimensional symmetric double, there is no proper nondegenerate invariant linear slice surjecting onto the original variables; any smaller realization must change the realization mechanism rather than delete such a slice.

\emph{Credit.} Credit recorded to ChatGPT (AI-assisted critique and derivation).

\phantomsection
\label{supp-note-05-026}
\paragraph{If an n-dimensional homogeneous map has a k-dimensional...}
\emph{Status.} Proof offered; independent review is not recorded.

If an n-dimensional homogeneous map has a k-dimensional partial-gradient block, it admits a weak symmetric dilation in dimension at most 2n-k; in particular every polynomial map has weak symmetric-dilation dimension at most 2n-1.

\emph{Credit.} Credit recorded to ChatGPT (AI-assisted critique and derivation).

\phantomsection
\label{supp-note-05-027}
\paragraph{For the displayed 19-dimensional block tensor...}
\emph{Status.} Recorded result or structural lead; it is not promoted beyond the evidence described here and in the project archive.

For the displayed 19-dimensional block tensor H(X, w, t)=(tQ(X)+t\textasciicircum{}2Bw, -q(X), 0), the weak symmetric-dilation invariant is at most 37-dim(w), hence at most 36.

\emph{Credit.} Credit recorded to ChatGPT (AI-assisted critique and derivation).

\phantomsection
\label{supp-note-05-028}
\paragraph{The admissible symmetric-overlap index kappa gives the lower bound...}
\emph{Status.} Proof offered; independent review is not recorded.

The admissible symmetric-overlap index kappa gives the lower bound sigma\_sym(H)>=2n-kappa\_ov(H); when the overlap form is nondegenerate, the partial-gradient construction gives the matching upper bound 2n-kappa\_nd(H).

\emph{Credit.} Credit recorded to ChatGPT (AI-assisted critique and derivation).

\phantomsection
\label{supp-note-05-029}
\paragraph{In the scalar-w block case, a 35-dimensional weak dilation extending...}
\emph{Status.} Recorded result or structural lead; it is not promoted beyond the evidence described here and in the project archive.

In the scalar-w block case, a 35-dimensional weak dilation extending the evident (w, t)-compression exists exactly when there is a nonzero vector a satisfying a\textasciicircum{}TB=0, a\textasciicircum{}TQ(X)=0, and D\_a q(X)=0 identically.

\emph{Credit.} Credit recorded to ChatGPT (AI-assisted critique and derivation).

\phantomsection
\label{supp-note-05-030}
\paragraph{For the principal sl2 conic on Sym\textasciicircum{}m(k\textasciicircum{}2), the integrability...}
\emph{Status.} Recorded result or structural lead; it is not promoted beyond the evidence described here and in the project archive.

For the principal sl2 conic on Sym\textasciicircum{}m(k\textasciicircum{}2), the integrability operator has determinant -m\textasciicircum{}\{m+2\}(m+2)\textasciicircum{}m; no direct sum of principal conics supports a noncollinear normalized collision.

\emph{Credit.} Credit recorded to ChatGPT (AI-assisted critique, derivation, and computation).

\phantomsection
\label{supp-note-05-031}
\paragraph{An everywhere-regular five-dimensional quadratic nilpotent line...}
\emph{Status.} Proof offered; independent review is not recorded.

An everywhere-regular five-dimensional quadratic nilpotent line pencil whose kernel map spans P4 has saturated filtration degrees (-4, -2, 0, 2, 4); every successive map is an isomorphism and the kernel map is the rational normal quartic.

\emph{Credit.} Credit recorded to ChatGPT (AI-assisted critique, derivation, and computation).

\phantomsection
\label{supp-note-05-032}
\paragraph{Every quadratic matrix annihilating the Veronese kernel k\_m factors...}
\emph{Status.} Proof offered; independent review is not recorded.

Every quadratic matrix annihilating the Veronese kernel k\_m factors uniquely as N=L\_mS\_m; iterated Hilbert-Burch descent preserves regular nilpotence, and a fully compatible canonical descent chain is necessarily a common scalar multiple of the principal chain.

\emph{Credit.} Credit recorded to ChatGPT (AI-assisted critique, derivation, and computation).

\phantomsection
\label{supp-note-05-033}
\paragraph{With the quartic kernel fixed at the principal point, the...}
\emph{Status.} Recorded result or structural lead; it is not promoted beyond the evidence described here and in the project archive.

With the quartic kernel fixed at the principal point, the nilpotent-pencil tangent space has dimension 16 and sl2 type V0+V2+V4+V6; fixing the descended cubic kernel leaves only scalar rescaling, so the other 15 directions are exactly relative descended-cubic motion.

\emph{Credit.} Credit recorded to ChatGPT (AI-assisted critique, derivation, and computation).

\phantomsection
\label{supp-note-05-034}
\paragraph{For a general five-dimensional quadratic pencil, integrable markings...}
\emph{Status.} Recorded result or structural lead; it is not promoted beyond the evidence described here and in the project archive.

For a general five-dimensional quadratic pencil, integrable markings lie on the proper discriminant det Phi\_N=0, and normalized collisions lie on the determinantal locus rank Theta\_N<=9 together with the open independence condition on the recovered marking vectors.

\emph{Credit.} Credit recorded to ChatGPT (AI-assisted critique, derivation, and computation).

\phantomsection
\label{supp-note-05-035}
\paragraph{For a global regular-type five-dimensional nilpotent quadratic...}
\emph{Status.} Recorded result or structural lead; it is not promoted beyond the evidence described here and in the project archive.

For a global regular-type five-dimensional nilpotent quadratic tensor with full-span collision-line kernel, the saturated rank-one defects are nested and have total codimension-two class 20H\textasciicircum{}2, combined Hilbert polynomial 10n\textasciicircum{}2+50n+81, and weighted effective increments 5Delta\_1+4Delta\_2+3Delta\_3+2Delta\_4+Delta\_5=20H\textasciicircum{}2.

\emph{Credit.} Credit recorded to ChatGPT (AI-assisted critique, derivation, and computation).

\subsection{Open problems and unresolved assertions}

\phantomsection
\label{supp-note-05-036}
\paragraph{What is the true smallest dimension of a cubic-homogeneous...}
\emph{Status.} Open problem or unresolved assertion.

What is the true smallest dimension of a cubic-homogeneous counterexample?

\phantomsection
\label{supp-note-05-037}
\paragraph{What is the exact minimal square-zero Druzkowski pairing rank for...}
\emph{Status.} Open problem or unresolved assertion.

What is the exact minimal square-zero Druzkowski pairing rank for the 19D cubic?

\subsection{Context, corrections, and evidence boundaries}

\phantomsection
\label{supp-note-05-038}
\paragraph{Claimed failure of the Mathieu conjecture for SU(3).}
\emph{Status.} Correction retained: the original positive claim is not adopted; only the displayed negative scope statement is current.

The cited dimension-varying implication does not establish that the Mathieu conjecture for the fixed group SU(3) is false.

\emph{Credit.} Credit recorded to Zihan Zhang (derivation, exposition).

\phantomsection
\label{supp-note-05-039}
\paragraph{Within a zero-diagonal five-by-five sparse matrix ansatz with at...}
\emph{Status.} Recorded result or structural lead; it is not promoted beyond the evidence described here and in the project archive.

Within a zero-diagonal five-by-five sparse matrix ansatz with at most four nonzero entries, no Keller map in the specified H\_C family realizes the normalized collision.

\phantomsection
\label{supp-note-05-040}
\paragraph{A second issue occurs in Zwart's fixed-highest-weight-vector proof.}
\emph{Status.} Reported context; not independently established here.

A second issue occurs in Zwart's fixed-highest-weight-vector proof.

\phantomsection
\label{supp-note-05-041}
\paragraph{No explicit Mathieu witness pair is currently recorded.}
\emph{Status.} Reported context; not independently established here.

The Mathieu argument currently proves existence of a violating pair, but no explicit closed-form pair of functions on SU(24) or SU(79) is recorded.  Making the generic Cartan-twist parameter explicit remains open.

\phantomsection
\label{supp-note-05-042}
\paragraph{The direct Euclidean 36-dimensional partial-gradient potential is...}
\emph{Status.} Corrected statement; the earlier stronger formulation is not adopted.

The direct Euclidean 36-dimensional partial-gradient potential is not Keller when B is nonzero; obtaining a 36-dimensional Keller or cover-preserving realization is a separate self-adjoint nilpotent matrix-completion problem.

\emph{Credit.} Credit recorded to ChatGPT (AI-assisted critique and derivation).

\endgroup
